%%%%%%%%%%%%%%%%%%%%%%%%%%%%%%%%%%%%%%%%%%%%%%%%%
\section{一圈阶的夸克准 PDF}
\label{sec:oneloop}

本节研究一圈阶的夸克准 PDF 及其重正化。费曼规范下，动量为 $p$ 的渐近夸克的夸克准 PDF 在一圈阶的费曼图见图~\ref{fig:oneloop}。图~\ref{fig:oneloop} 中的图 (a) 给出
\begin{align}
\begin{split}\label{eq:1a0}
M_{1a}=&\frac{e^{i p_z \xi_z}}{p_z}\frac{1}{N_c} \text{Tr}_c[T^aT^a] \int_a^{\xi_z-2a}dr_1\int_{r_1+a}^{\xi_z-a}dr_2\\
&\times \int \frac{d^4l}{(2\pi)^4} \,e^{-i p_z \xi_z} \,e^{il_z(r_2-r_1)} \left(\frac{-ig_{\mu\nu}}{l^2}\right) \\
&\times (-ig_s n_z^\mu)(-ig_s n_z^\nu)  \text{Tr}\left[\frac{1}{2}\, \slashed{p}\, \frac{1}{2}\gamma_z\right]   \\
=& \frac{\alpha_s C_F}{4 i \pi^3} \int_a^{\xi_z-2a}dr_1\int_{r_1+a}^{\xi_z-a}dr_2 \int d^4l\frac{e^{il_z(r_2-r_1)}}{l^2}
\end{split}
\end{align}
其中我们在沿规范链接的两场之间引入了截断``$a$''，用以正规化潜在的线性紫外发散\footnote{尽管紫外发散的显式结果依赖于紫外正规子，但我们的结论（如幂次计数规则与重正化结构）并不依赖于它们。}。式 (9) 中 $r_1$ 积分的上限为 $\xi_z-2a$，因为它需要与 $r_2$ 积分的上限保持间距 $a$。我们将证明这个截断足以正规化此图中的所有紫外发散，因此此处不引入 DR。为后续计算方便，引入矢量 $\bar{l}^\mu$，它与 $l^\mu$ 相同但没有 $z$ 方向分量：$l^\mu = \bar{l}^\mu + l_z n_z^\mu$，$l^2=\bar{l}^2-l_z^2$。我们把对 $\bar{l}^\mu$ 的积分称为三维（3-D）积分，把对 $l^\mu$ 的积分称为四维（4-D）积分。记 $d^4l = d^{3}\bar{l}\, dl_z$，考虑如下相空间积分：
\begin{eqnarray}\label{eq:3dexpansion}
\int \frac{d^{3}\bar{l}}{l^2} &=&
\int \frac{d^{3}\bar{l}}{\bar{l}^2-l_z^2}
\nonumber\\
&=& \int d^{3}\bar{l} \left( \frac{1}{\bar{l}^2}+ \frac{l_z^2}{(\bar{l}^2-l_z^2)\bar{l}^2}\right).
\end{eqnarray}
式~(\ref{eq:3dexpansion}) 的第一项线性发散，但其系数正比于 $\int dl_z e^{il_z(r_2-r_1)} = 2\pi \delta(r_2-r_1)$，由于 $r_2$ 与 $r_1$ 不能处于同一点，该项为零。因此，上述三维积分实际上是有限的。另一方面，容易看出，只要 $\bar{l}^\mu$ 保持有限，即使 $a\to 0$，$l_z$ 的积分也是有限的。所以，式~\eqref{eq:1a0} 只可能在 $l^\mu$ 的所有分量都趋于无穷的区域出现紫外发散。结果是，间距``$a$''足以正规化所有紫外发散，给出
\begin{align}\label{eq:1a}
\begin{split}
M_{1a}\overset{\text{div}}{=}-\frac{\alpha_s C_F}{\pi}\frac{|\xi_z|}{a}+\frac{\alpha_s C_F}{\pi}\ln \frac{|\xi_z|}{a},
\end{split}
\end{align}
其中我们已包含 $\xi_z < 0$ 的情形。


%%%%%%%%%%%%%%%%%%%%%%%%%%%%%%%
\begin{figure}[t]
\begin{center}
\includegraphics[width=2.5in]{oneloop}
\caption{\label{fig:oneloop}
动量为 $p$ 的渐近夸克的夸克准 PDF 在一圈阶的费曼图。}
\end{center}
\end{figure}
%%%%%%%%%%%%%%%%%%%%%%%%%%%%%%%


由图~\ref{fig:oneloop}，图 (b) 的贡献为
\begin{align}
\begin{split}\label{eq:1b0}
M_{1b}=& g_s^2C_F \int_a^{\xi_z-a}dr\int \frac{d^4l}{(2\pi)^4}\,
\frac{e^{il_z r}}{l^2(p-l)^2}
 \\
& \times \frac{1}{p_z} \text{Tr}\left[\frac{1}{2}\, \slashed{p}\, \frac{1}{2}\gamma_z(\slashed{p}-\slashed{l})\gamma_z\right],
\end{split}
\end{align}
这里再次使用间距``$a$''作为正规子。对于三维积分，正比于 $\slashed{l}$ 的项有潜在的对数发散：
\begin{align}
\int \frac{l^\mu\,d^{3}\bar{l}}{l^2(p-l)^2}.
\end{align}
然而，上述积分是有限的：因为要提取潜在的对数发散，可以在积分中令外部物理标度 $p\to 0$，此时它成为 $\bar{l}^\mu$ 的奇函数，积分为零。这就解释了间距 $a$ 为什么能够正规化式~\eqref{eq:1b0} 的紫外发散。图 (b) 的紫外发散为
\begin{align}\label{eq:1b}
\begin{split}
M_{1b}\overset{\text{div}}{=}-\frac{\alpha_s C_F}{2\pi} \ln \frac{|\xi_z|}{a},
\end{split}
\end{align}
其中我们同样包含了 $\xi_z<0$ 的情形。

类似地，由图~\ref{fig:oneloop} 中的图 (c) 可得
\begin{align}\label{eq:1c0}
\begin{split}
M_{1c}=&\frac{i g_s^2 \mu_r^{2\epsilon}}{2p_z}(1-\epsilon)C_F\int\frac{d^dl}{(2\pi)^d} \frac{\text{Tr}[\slashed{p}\gamma_z {\slashed{p}}(\slashed{p}-\slashed{l})]}{p^2l^2(p-l)^2}\\
=&i g_s^2 \mu_r^{2\epsilon} C_F(1-\epsilon)\int\frac{d^dl}{(2\pi)^d} \frac{1}{l^2(p-l)^2},
\end{split}
\end{align}
其中时空维数定义为 $d=4-2\epsilon$。显然，当 $p^2=0$ 时，式~(\ref{eq:1c0}) 中的积分在 DR 下为零，原因是 $1/\epsilon$ 的紫外极点与红外极点表观相消。只有当 $l^\mu$ 的所有分量都趋于无穷时，此图的紫外发散才是对数型的，在 DR 中为
\begin{align}\label{eq:1c}
\begin{split}
M_{1c}\overset{\text{div}}{=}-\frac{\alpha_s C_F}{4\pi} \frac{1}{\epsilon}.
\end{split}
\end{align}
众所周知，在这一阶及更高阶，无质量夸克自能图的紫外发散可以通过夸克场重正化消除。

图~\ref{fig:oneloop} 中的图 (d) 给出
\begin{align}
\begin{split}\label{eq:1d0}
M_{1d}=& -i g_s^2C_F e^{i p_z\xi_z} \int \frac{d^4k}{(2\pi)^4} e^{-ik_z \xi_z} \frac{1}{k^4 (p-k)^2} \\
&\times \frac{1}{4p_z } \text{Tr}[\slashed{p}\gamma_\mu\slashed{k}\gamma_z\slashed{k}\gamma^\mu]
\end{split}
\end{align}
其中 $k=p-l$。
由于振荡因子 $e^{-ik_z \xi_z}$，来自 $k_z$ 趋于无穷区域的贡献受到强烈压低。因此，它在有限 $\xi_z$ 下是紫外有限的，尽管当 $\xi_z \to 0$ 时发散：
\begin{align}\label{eq:1d}
\begin{split}
M_{1d}\overset{\text{div}}{=}-\frac{\alpha_s C_F}{2\pi} \ln(|\xi_zp_z|).
\end{split}
\end{align}

综上，动量为 $p$ 的夸克的夸克准 PDF 在任意 $\xi_z$ 处紫外发散的总一圈贡献为
\begin{align}\label{eq:1}
\begin{split}
M^{(1)}&\overset{\text{div}}{=}M_{1a}+2\times M_{1b}+ 2\times\frac{1}{2}M_{1c}+M_{1d}\\
&=\frac{\alpha_s C_F}{\pi}\left(-\frac{|\xi_z|}{a} - \frac{1}{4\epsilon} - \frac{1}{2} \ln(|\xi_zp_z|) \right),
\end{split}
\end{align}
其中 $\ln(|\xi_zp_z|)$ 项并非紫外发散，而是在 $\xi_z \to 0$ 时发散
\footnote{如果取 $\xi_z\to 0$，我们所计算的便是对局域矢流的一圈修正。式~\eqref{eq:1} 中有限的 $\xi_z$ 实际上正规化了圈顶点图（图~\ref{fig:oneloop}(d)）的紫外发散，而图~\ref{fig:oneloop}(c) 自能图的紫外发散由维数正规化处理。从图~\ref{fig:oneloop}(d) 出发，我们在固定 $\xi_z$ 的条件下取 $\epsilon\to 0$ 得到了 $\ln(|\xi_z p_z|)$ 项。如果需要取 $\xi_z\to 0$ 以模仿局域流，则应保持 $\epsilon$ 有限以正规化一圈紫外发散，这将使 $\ln(|\xi_z p_z|)$ 变为 $- \frac{1}{2\epsilon}$。于是从式~\eqref{eq:1} 我们发现：在固定 $\epsilon$ 与 $a$ 的条件下取 $\xi_z\to 0$ 时，矢流的紫外发散在一圈水平消失，这与流守恒的期望一致。}。
我们发现，在这一阶，紫外发散只来自所有圈动量都趋于无穷的区域，因此紫外发散在坐标空间中是局域的。下一节将证明，这一性质在 QCD 微扰论所有阶都保持成立。


%%%%%%%%%%%%%%%%%%%%%%%%%%%%%%%
\begin{figure}[ht]
\begin{center}
\includegraphics[width=3.3in]{oneloopg}
\caption{\label{fig:oneloopg}
动量为 $p$ 的渐近胶子的夸克准 PDF 在一圈阶的费曼图。}
\end{center}
\end{figure}
%%%%%%%%%%%%%%%%%%%%%%%%%%%%%%%

费曼规范下，动量为 $p$ 的渐近胶子的夸克准 PDF 在一圈阶的费曼图见图~\ref{fig:oneloopg}，外加图 (b) 的复共轭图。下一节的一般性论证将表明，准 PDF 所有图的紫外发散只能来自 4-D 积分，因而在坐标空间局域。如果保持 $\xi_z$ 有限，图~\ref{fig:oneloopg} 的所有一圈图都不能是局域的，正如~\ref{fig:oneloop} 中的图 (d) 一样，因此图~\ref{fig:oneloopg} 中的这些一圈图必定是紫外有限的。但是，当 $\xi_z\to 0$ 时它们可以发散。为了演示这一特征，以图~\ref{fig:oneloopg} 中的图 (a) 为例。图~\ref{fig:oneloopg}(a) 给出
\begin{align}
\begin{split}\label{eq:2a0}
M_{2a}\propto& \int_0^{\xi_z} dr_1 \int_{r_1}^{\xi_z} dr_2 \int d^4l\, e^{-il_z \xi_z} \frac{l_z}{l^2} \\
=&\frac{\xi_z^2}{2}\int dl_z \, e^{-il_z \xi_z} \, l_z \int d^{3}\bar{l} \left( \frac{1}{\bar{l}^2}+ \frac{l_z^2}{(\bar{l}^2-l_z^2)\bar{l}^2}\right),
\end{split}
\end{align}
它看起来在三维积分下线性紫外发散。然而，与图~\ref{fig:oneloop} 中图 (a) 的情形类似，线性发散项正比于 $\delta^\prime(\xi_z)$，对有限的 $\xi_z$ 为零。第二项在对 $\bar{l}$ 积分后有限，给出
\begin{align}
\begin{split}\label{eq:2a}
& \frac{\xi_z^2}{2}\int dl_z \, e^{-il_z \xi_z} \, l_z \int d^{3}\bar{l} \frac{l_z^2}{(\bar{l}^2-l_z^2)\bar{l}^2}\\
\propto&\frac{\xi_z^2}{2}\int dl_z \, e^{-il_z \xi_z} \, \frac{l_z^3}{|l_z|}\\
=&\frac{2i}{\xi_z},
\end{split}
\end{align}
其中的比例关系可以由量纲计数简单地得到。因此我们发现，图~\ref{fig:oneloopg}(a) 是紫外有限的，但在 $\xi_z \to 0$ 时表现为 $1/\xi_z$。类似地，我们发现图~\ref{fig:oneloopg}(b,c) 也是紫外有限的，这表明在一圈水平上，夸克准 PDF 在紫外重正化下不会从胶子获得混合。


继续之前需要指出，上面发现的紫外行为与 PDF 的显著不同。PDF 的紫外发散来自圈动量 $(l_+, l_-, \vec{l}_\perp)\sim(1,\lambda^2,\lambda)$（$\lambda\to\infty$）的区域。于是，由于动量分量 $l_+\sim {\cal O}(1)$，普通 PDF 的紫外发散在坐标空间的``$-$''方向上是非局域的。正是这一事实使得 PDF 的重正化是卷积形式、与沿``$-$''方向的带算符的所有扭度2 PDF 相混合，而如我们将证明的，准 PDF 的重正化是一个乘法因子。

我们进一步指出，即便紫外发散（$1/a$、$\ln(a)$ 或 $1/\epsilon$）已被重正化，一圈图的 $\ln(|\xi_z|)$ 依赖（见式~\ref{eq:1d} 为例）与 $1/\xi_z$ 依赖（见式~\ref{eq:2a} 为例）仍预示着准 PDF 在 $\xi_z\to 0$ 时发散；如上一节所述，这导致若简单地对坐标空间准 PDF 作傅里叶变换，就难以得到良定的动量空间准 PDF。这一行为在 QCD 微扰论的更高阶必然仍然成立。
